\chapter{Further Properties of Brownian Motion}

\section{Framework}

Let $(B_t)_{t\ge0}$ be a one--dimensional standard Brownian motion on a probability space
$(\Omega,\mathcal F,\Pp)$.
We work with the completed, right--continuous natural filtration
\[
\widehat{\mathcal F}_t
:= \sigma\{B_s:0\le s\le t\}\vee\mathcal N,
\]
where $\mathcal N$ denotes the $\Pp$--null sets.

A random time $T:\Omega\to[0,\infty]$ is called a \emph{stopping time} if
\[
\{T\le t\}\in\widehat{\mathcal F}_t
\quad\text{for all }t\ge0.
\]

\section{Exit Times and the Skorokhod Embedding}

Imagine a factory producing a sensitive electronic component. One important quantity — say voltage deviation — fluctuates randomly over time due to usual degradation. The deviation has the folloing properties: starts at $0V$, it fluctuates unpredictably over time and we cannot influence the process once it starts. We assume that the fluctuation is well modeled by Brownian motion. We want the final recorded deviation to satisfy a precise statistical requirement: with probability 50\%, the deviation should be within $+1V$,
and with probability 50\%, within $-1V$ unit. In other words, we want to identify a situation where 50\% of our component degrades beyond our tolerance. Skorokhod embedding will give us a stopping-time rule that, when applied repeatedly to identical components, produces a prescribed distribution of final outcomes across the population. Actually we will see by stopping the process at a carefully chosen random time, we can \emph{manufacture} a prescribed probability distribution.


\subsection{Exit time from an interval}
We will need a class of stopping times, namely, \emph{hitting/exit times}.
If we start at $0$ and look at the first time the path leaves a bounded interval $(a,b)$ with $a<0<b$,
then the stopped value $B_\gamma$ can only be the endpoint $a$ or $b$.
So although $B_\gamma$ is simple (a two--point distribution), it is already \emph{tunable}:
\begin{itemize}
\item the probability of exiting at $a$ versus $b$ depends on the geometry of the interval,
\item and the expected exit time $\E[\gamma]$ can be computed explicitly.
\end{itemize}

Fix $a<0<b$ and define the exit time 
\[
\gamma := \inf\{t>0:\, B_t\notin(a,b)\}.
\]

\begin{proposition}
The following statements hold:
\begin{enumerate}
\label{prop:exit_time_prop}
\item $\Pp(\gamma<\infty)=1$.
\item $\Pp(B_\gamma=a)=\frac{b}{\,b-a\,},
\qquad
\Pp(B_\gamma=b)=\frac{-a}{\,b-a\,}.$
\item $\E[\gamma]=-ab$.
\end{enumerate}

\end{proposition}

We will need the following thoerem to proof propostion \ref{prop:exit_time_prop}.

\begin{theorem}[Optional Stopping Theorem]
Let $(M_t)_{t\ge0}$ be a martingale and $T$ a stopping time such that one of the following
conditions holds:
\begin{itemize}
\item $T$ is almost surely bounded, or
\item $(M_{t\wedge T})_{t\ge0}$ is uniformly integrable.
\end{itemize}
Then
\[
\E[M_T]=\E[M_0].
\]
\end{theorem}

\begin{proof}[Proof of Proposition \ref{prop:exit_time_prop}]
We prove each statement separately. The optional stopping theorem (without proof) is the main tool for statement 2 and 3. 

\Step{1: Almost sure finiteness of $\gamma$}

Since Brownian motion is recurrent in one dimension, it exits every bounded interval almost surely. Hence
\[
\Pp(\gamma<\infty)=1.
\]

\medskip

\Step{2: Exit probabilities}

For $\gamma_n:=\gamma\wedge n$, optional stopping applied to the martingale $(B_t)$ yields
\[
\E[B_{\gamma_n}]=0.
\]
Letting $n\to\infty$ gives $\E[B_\gamma]=0$. Since $B_\gamma\in\{a,b\}$,
\[
a\,\Pp(B_\gamma=a)+b\,\Pp(B_\gamma=b)=0,
\]
which together with normalization yields the stated probabilities.

\Step{3: Expectation of $\gamma$}

Applying optional stopping to the martingale $(B_t^2-t)$ gives
\[
\E[B_\gamma^2]=\E[\gamma].
\]
Since $B_\gamma\in\{a,b\}$, this yields $\E[\gamma]=-ab$.
\end{proof}

\subsection{Skorokhod embedding}

\paragraph{The embedding idea.}
The Skorokhod embedding problem asks:
given a target law $F$ on $\R$ with mean $0$ and finite variance, can we find a stopping time $T$ such that
$B_T\sim F$?
The guiding idea is to build $T$ by \emph{randomizing the interval}.
Instead of a deterministic $(a,b)$, we choose random endpoints $(\alpha,\beta)$ with $\alpha<0<\beta$ and then stop
\[
T=\inf\{t>0:\,B_t\notin(\alpha,\beta)\}.
\]
Conditionally on $(\alpha,\beta)=(a,b)$ we already know the law of $B_T$ (it is concentrated on $\{a,b\}$ with explicit weights),
so the unconditional law of $B_T$ becomes a \emph{mixture} (integral) over $(a,b)$.
By choosing the mixing measure $\mu$ appropriately, the mixture can be made equal to the prescribed $F$.

\medskip

Let $F$ be a probability measure on $\R$ with
\[
\int_\R x\,F(dx)=0,
\qquad
\int_\R x^2\,F(dx)=\sigma^2<\infty.
\]

Split $F$ into its positive and negative parts,
\[
F^+ := F|_{[0,\infty)},\qquad F^- := F|_{(-\infty,0)}.
\]

One can construct a probability measure $\mu$ on
$(-\infty,0)\times(0,\infty)$ such that if
$(\alpha,\beta)\sim\mu$ is independent of $B$ and
\[
T:=\inf\{t>0:\, B_t\notin(\alpha,\beta)\},
\]
then
\[
B_T\sim F
\quad\text{and}\quad
\E[T]=\sigma^2.
\]

\begin{theorem}[Skorokhod embedding]
There exists a stopping time $T$ such that
\[
B_T\sim F
\quad\text{and}\quad
\E[T]=\sigma^2.
\]
\end{theorem}

\begin{proof}
Define the normalizing constant $\gamma>0$ by
\[
\gamma^{-1}
:= \int_{0}^{\infty} b\,F(db)
= -\int_{-\infty}^{0} a\,F(da)
= \frac12\int_{\R} |x|\,F(dx),
\]
where the equalities follow from $\int_{\R} x\,F(dx)=0$.

On $(-\infty,0)\times(0,\infty)$ define the probability measure
\[
\mu(da,db)
:= \gamma\,\frac{1}{b-a}\,F^+(db)\,F^-(da).
\]
(That $\mu$ is a probability measure is exactly the normalization encoded by $\gamma^{-1}$.)

Let $(\alpha,\beta)\sim\mu$ be independent of the Brownian motion $B$, and set
\[
T:=\inf\{t>0:\,B_t\notin(\alpha,\beta)\}.
\]
Conditionally on $\{(\alpha,\beta)=(a,b)\}$, Proposition \ref{prop:exit_time_prop} gives
\[
\Pp(B_T\in db\mid \alpha=a,\beta=b)=\frac{-a}{b-a}\,\delta_b(db),
\qquad
\Pp(B_T\in da\mid \alpha=a,\beta=b)=\frac{b}{b-a}\,\delta_a(da),
\]
where $\delta_x$ denotes the Dirac measure at $x$.

\medskip
\LabelProof{Law of $B_T$}
Fix a Borel set $E\subset(0,\infty)$. Using conditioning and the definition of $\mu$,
\begin{align*}
\Pp(B_T\in E)
&=\int_{-\infty}^{0}\int_{0}^{\infty}
\Pp(B_T\in E\mid \alpha=a,\beta=b)\,\mu(da,db)\\
&=\int_{-\infty}^{0}\int_{0}^{\infty}
\left(\frac{-a}{b-a}\mathbf{1}_E(b)\right)\,
\gamma\frac{1}{b-a}\,F^+(db)\,F^-(da)\\
&=\int_{0}^{\infty}\mathbf{1}_E(b)\,
\left[\gamma\int_{-\infty}^{0}\frac{-a}{(b-a)^2}\,F^-(da)\right]F^+(db).
\end{align*}
By construction the notes, the bracketed factor is equal to $1$ for $F^+$--a.e.\ $b$,
so the last line equals $F^+(E)$, i.e.
\[
\Pp(B_T\in E)=F^+(E).
\]
Analogously, for Borel $E\subset(-\infty,0)$ one obtains
\[
\Pp(B_T\in E)=F^-(E).
\]
Since $F=F^-+F^+$, it follows that $B_T\sim F$.

\medskip
\LabelProof{Expectation of $T$}
By Proposition~2.1, conditionally on $(\alpha,\beta)=(a,b)$ we have
\[
\E[T\mid \alpha=a,\beta=b]= -ab.
\]
Therefore,
\begin{align*}
\E[T]
&=\int_{-\infty}^{0}\int_{0}^{\infty} (-ab)\,\mu(da,db)\\
&=\int_{-\infty}^{0}\int_{0}^{\infty} (-ab)\,
\gamma\frac{1}{b-a}\,F^+(db)\,F^-(da).
\end{align*}
Rewriting $-ab/(b-a)$ in the form used in the notes and applying $\int x\,F(dx)=0$
yields the identity
\[
\E[T]=\int_{\R} x^2\,F(dx)=\sigma^2.
\]
This completes the proof.
\end{proof}

\begin{remark}
The construction relies on the fact that Brownian motion exits intervals only at the
endpoints and that both the exit distribution and the expected exit time are explicitly
known.
\end{remark}

\begin{remark}
\paragraph{Why the variance appears as $\E[T]$?}
The second requirement $\E[T]=\sigma^2$ is not accidental.
Brownian motion satisfies $\E[B_t^2]=t$, so time plays the role of variance.
Once we enforce $B_T\sim F$, we automatically have $\E[B_T^2]=\sigma^2$.
The proof then uses the explicit formula $\E[T\mid \alpha,\beta]=-\,\alpha\beta$ from Proposition \ref{prop:exit_time_prop} and integrates it against $\mu$:
the construction of $\mu$ is designed so that this integral equals $\sigma^2$.
In this way, the embedding identifies the ``amount of time'' needed for Brownian motion to accumulate the required second moment.    
\end{remark}


\section{Brownian Motion as a Markov Process}

% --- Intuitive introduction to Section 3: Brownian Motion as a Markov Process ---

So far, Brownian motion has been treated primarily as a stochastic process with explicit
finite--dimensional distributions.
In this section we change perspective and view Brownian motion as a \emph{Markov process}.
The central theme is that the future evolution of the path depends on the past only through the
current position, and that this fact can be expressed analytically through transition operators,
generators, and martingales.

Understanding Brownian motion as a Markov process prepares the ground for the strong Markov property
and for path transformations such as reflection, which rely crucially on this viewpoint.


\subsection{Markov property at deterministic times}

\paragraph{Idea.}
The Markov property is not an additional assumption for Brownian motion;
it is a direct consequence of independent increments.
Writing
\[
B_{t+s}=B_s+(B_{t+s}-B_s),
\]
the increment $B_{t+s}-B_s$ is independent of everything that happened before time $s$ and has the same
law as $B_t$.
This immediately implies that conditional expectations of the form
$\E[f(B_{t+s})\mid\widehat{\mathcal F}_s]$ can only depend on $B_s$.

\medskip

Let $(B_t)_{t\ge0}$ be a standard Brownian motion with respect to the completed,
right--continuous filtration $(\widehat{\mathcal F}_t)_{t\ge0}$.

Fix $s,t\ge0$. We decompose
\[
B_{t+s} = B_s + (B_{t+s}-B_s).
\]
By the independent increments property, the increment $B_{t+s}-B_s$ is independent of
$\widehat{\mathcal F}_s$ and has the same law as $B_t$.

Let $f:\R\to\R$ be bounded and measurable. Then
\begin{align*}
\E\!\left[f(B_{t+s}) \mid \widehat{\mathcal F}_s\right]
&= \E\!\left[f\bigl(B_s + (B_{t+s}-B_s)\bigr) \mid \widehat{\mathcal F}_s\right] \\
&= \E\!\left[f(B_s + B_t') \mid \widehat{\mathcal F}_s\right],
\end{align*}
where $B_t'$ is an independent copy of $B_t$.

Since $B_t'$ is independent of $\widehat{\mathcal F}_s$, the conditional expectation reduces to a
function of $B_s$ only. Hence there exists a measurable function $g:\R\to\R$ such that
\[
\E\!\left[f(B_{t+s}) \mid \widehat{\mathcal F}_s\right] = g(B_s),
\]
where
\[
g(x) := \E[f(x+B_t)].
\]

This shows that Brownian motion satisfies the (time--homogeneous) Markov property.

\subsection{The Markov family \texorpdfstring{$(\mathbb P_x)_{x\in\R}$}{}}
Before analyzing the evolution of a stochastic process in time, we introduce the concept of a Markov family, which formalizes the idea that the probabilistic behavior of the process can be restarted from its current state. Roughly speaking, once the present state of the system is known, the future evolution is governed by the same probabilistic laws, independently of how the system arrived at that state.

Returning to the manufacturing example, this corresponds to the following intuition: if the degradation process of a component is temporarily suspended (for instance, by turning the component off) and later resumed, then—given the current level of degradation—we expect the subsequent evolution to follow the same statistical behavior as if the process had just started from that level. The past history of the component matters only through its present state, not through the specific path taken to reach it.

\medskip 

For each $x\in\R$, define a probability measure $\mathbb P_x$ by
\[
\mathbb P_x(A) := \mathbb P\bigl((x+B_t)_{t\ge0}\in A\bigr),
\qquad A\in\mathcal F.
\]
Under $\mathbb P_x$, the process starts at $B_0=x$ and has the same increments as standard Brownian
motion.

The family
\[
\bigl(\Omega,\mathcal F,(\widehat{\mathcal F}_t)_{t\ge0},(B_t)_{t\ge0},(\mathbb P_x)_{x\in\R}\bigr)
\]
is therefore a time--homogeneous Markov process with state space $\R$.

\paragraph{Why introduce the family $(\mathbb P_x)_{x\in\R}$?}
Once we realize that only the current position matters, it becomes natural to consider Brownian motion
starting from an arbitrary point $x$.
The family of laws $(\mathbb P_x)_{x\in\R}$ encodes this idea:
under $\mathbb P_x$ the process starts at $B_0=x$ and then evolves like standard Brownian motion.
This viewpoint allows us to treat Brownian motion as a time--homogeneous Markov process with state space $\R$,
rather than as a single process started at $0$.



\subsection{Transition semigroup}

Let us now investigate what happens when we enforce the Markov property
\begin{equation}\label{eq:Markov_conditional}
\mathbb{E}_x\!\left[f(X_{t+s}) \,\middle|\, \mathcal{F}_s\right]
=
\mathbb{E}_{X_s}\!\left[f(X_t)\right]
\qquad \mathbb{P}_x\text{-a.s.} \; \forall s,t\ge 0.
\end{equation}
Both sides of \eqref{eq:Markov_conditional} are \emph{random variables}:
the left-hand side is $\mathcal{F}_s$--measurable by definition of conditional expectation,
and the right-hand side depends on $\omega$ only through the random state $X_s(\omega)$. We can have another view on this equation if we define an operator $(P_t)_{t\ge 0}$ by
\[
(P_t f)(x) := \mathbb{E}_x\!\left[f(X_t)\right], \qquad x\in E.
\]
With this notation,
\[
\mathbb{E}_{X_s}\!\left[f(X_t)\right] = (P_t f)(X_s),
\]
i.e.\ the right-hand side of \eqref{eq:Markov_conditional} is the function $P_t f$
evaluated at the random point $X_s$.

\paragraph{Averaging both sides.}For the left-hand side we use
the tower property, hence
\[
\mathbb{E}_x\!\left[\mathbb{E}_x\!\left[f(X_{t+s}) \,\middle|\, \mathcal{F}_s\right]\right]
=
\mathbb{E}_x\!\left[f(X_{t+s})\right].
\]
For the right-hand side we use our previous notation
\[
\mathbb{E}_x\!\left[\mathbb{E}_{X_s}\!\left[f(X_t)\right]\right]
=
\mathbb{E}_x\!\left[(P_t f)(X_s)\right].
\]
Putting these together yields
\begin{equation}\label{eq:semigroup_intermediate}
\mathbb{E}_x\!\left[f(X_{t+s})\right]
=
\mathbb{E}_x\!\left[(P_t f)(X_s)\right].
\end{equation}

But by definition of $P_{t+s}$, the left-hand side of \eqref{eq:semigroup_intermediate} equals
\[
\mathbb{E}_x[f(X_{t+s})] = (P_{t+s} f)(x).
\]
By definition of $P_s$ applied to the function $P_t f$, the right-hand side equals
\[
\mathbb{E}_x[(P_t f)(X_s)] = (P_s(P_t f))(x).
\]
Therefore,
\[
(P_{t+s} f)(x) = (P_s(P_t f))(x), \qquad \forall\,x\in E,\ \forall\,f.
\]
Equivalently, as an operator identity,
\[
P_{t+s} = P_s \circ P_t, \qquad s,t\ge 0,
\]
which is precisely the \emph{semigroup property} and $P_t$ is what we call the transition semigroup (operator).

\paragraph{Interpretation.}
The transition semigroup $(T_t)_{t\ge0}$ describes how functions of the state are transported forward in time by
\[
T_t f(x)=\E_x[f(B_t)].
\]
Analytically, $T_t$ acts by convolution with the Gaussian kernel.
Conceptually, $T_t f(x)$ is the expected value of $f$ after the system has evolved for time $t$
starting from $x$.
The semigroup property $T_{t+s}=T_tT_s$ is simply the probabilistic restatement of the fact that Brownian motion
can be run in two successive time intervals.

\medskip

\begin{definition}[Transition Semigroup]
Define, for $f\in C_b(\R)$ and $t\ge0$,
\[
T_t f(x) := \E_x[f(B_t)].
\]

By the Gaussian transition density of Brownian motion,
\[
T_t f(x)
=
\int_{\R} f(y)\,p_t(x,y)\,dy,
\qquad
p_t(x,y)
=
\frac{1}{\sqrt{2\pi t}}
\exp\!\left(-\frac{(x-y)^2}{2t}\right).
\]

The family $(T_t)_{t\ge0}$ satisfies:
\begin{itemize}
\item $T_0 f = f$,
\item $T_{t+s} = T_t T_s$ (semigroup property),
\item $\|T_t f\|_\infty \le \|f\|_\infty$.
\end{itemize}

$T_t$ is the Transistion Semigroup.
    
\end{definition}


\subsection{Generator of Brownian motion}
The transition semigroup associated with a Markov process describes how the law of the
process evolves over finite time intervals. While this global description is sufficient
to propagate expectations forward in time, it does not reveal the local mechanism that
drives the evolution of the process. In many situations, one is interested in
understanding how the process changes instantaneously, rather than after a fixed positive
time.

The generator captures this infinitesimal behavior. It measures the first--order rate of
change of expected values at time zero and thus encodes the local dynamics that produce
the semigroup. In this sense, the generator plays the same role as a differential operator
in the theory of deterministic dynamical systems: it provides a concise and structural
description of the evolution, from which the entire semigroup can be recovered.

\begin{definition}
For functions $f$ for which the limit exists, define the generator
\[
Gf(x) := \lim_{h\downarrow0} \frac{T_h f(x)-f(x)}{h}.
\]    
\end{definition}


\begin{proposition}[Proposition 3.1]
If $f\in C_b^2(\R)$, then
\[
Gf(x)=\frac12 f''(x).
\]
\end{proposition}

\begin{proof}
Fix $x\in\R$ and $h>0$. By definition of the transition semigroup and the Gaussian kernel,
\begin{align*}
T_h f(x)
&=\E_x[f(B_h)]
=\int_{\R} f(y)\,p_h(x,y)\,dy\\
&=\int_{\R} f(y)\,\frac{1}{\sqrt{2\pi h}}
\exp\!\left(-\frac{(y-x)^2}{2h}\right)\,dy.
\end{align*}
Make the change of variables $y=x+z$ (so $dy=dz$). Then
\begin{equation}\label{eq:Th_change_var}
T_h f(x)
=\int_{\R} f(x+z)\,\frac{1}{\sqrt{2\pi h}}
\exp\!\left(-\frac{z^2}{2h}\right)\,dz.
\end{equation}
Equivalently, if $Z\sim\mathcal N(0,h)$, then \eqref{eq:Th_change_var} is $T_h f(x)=\E[f(x+Z)]$.

Now apply Taylor's theorem at $x$. For each $z\in\R$ there exists $\theta=\theta(z)\in(0,1)$ such that
\begin{equation}\label{eq:taylor}
f(x+z)
=
f(x)+f'(x)z+\frac12 f''(x)z^2
+
\frac12\bigl(f''(x+\theta z)-f''(x)\bigr)z^2.
\end{equation}
Insert \eqref{eq:taylor} into \eqref{eq:Th_change_var} and subtract $f(x)$:
\begin{align*}
T_h f(x)-f(x)
&=
f'(x)\int_{\R} z\,\frac{1}{\sqrt{2\pi h}}
e^{-z^2/(2h)}\,dz\\
&\quad+\frac12 f''(x)\int_{\R} z^2\,\frac{1}{\sqrt{2\pi h}}
e^{-z^2/(2h)}\,dz\\
&\quad+\frac12 \int_{\R}\bigl(f''(x+\theta z)-f''(x)\bigr)z^2\,
\frac{1}{\sqrt{2\pi h}}e^{-z^2/(2h)}\,dz.
\end{align*}

\medskip
\Step{1: compute the Gaussian moments}
The first integral is $0$ because the integrand is odd:
\[
\int_{\R} z\,\frac{1}{\sqrt{2\pi h}}e^{-z^2/(2h)}\,dz=0.
\]
The second integral is the variance of $\mathcal N(0,h)$:
\[
\int_{\R} z^2\,\frac{1}{\sqrt{2\pi h}}e^{-z^2/(2h)}\,dz = h.
\]
Hence
\begin{equation}\label{eq:main_term}
T_h f(x)-f(x)=\frac12 f''(x)\,h + R_h(x),
\end{equation}
where the remainder term is
\[
R_h(x)
:=\frac12 \int_{\R}\bigl(f''(x+\theta z)-f''(x)\bigr)z^2\,
\frac{1}{\sqrt{2\pi h}}e^{-z^2/(2h)}\,dz.
\]

\medskip
\Step{2: show $R_h(x)=o(h)$ as $h\downarrow0$}
Fix $\varepsilon>0$. Since $f''$ is continuous at $x$, there exists $\delta>0$ such that
\[
|u-x|\le \delta \quad\Longrightarrow\quad |f''(u)-f''(x)|\le \varepsilon.
\]
Split the integral defining $R_h(x)$ into $|z|\le\delta$ and $|z|>\delta$.

\emph{(i) Region $|z|\le\delta$.}
For $|z|\le\delta$ we have $|x+\theta z-x|\le |z|\le\delta$, hence
$|f''(x+\theta z)-f''(x)|\le\varepsilon$ and thus
\[
\left|\frac12\int_{|z|\le\delta}\!\!\bigl(f''(x+\theta z)-f''(x)\bigr)z^2\,
\frac{1}{\sqrt{2\pi h}}e^{-z^2/(2h)}\,dz\right|
\le \frac{\varepsilon}{2}\int_{\R} z^2\,\frac{1}{\sqrt{2\pi h}}e^{-z^2/(2h)}\,dz
=\frac{\varepsilon}{2}\,h.
\]

\emph{(ii) Region $|z|>\delta$.}
Since $f''$ is bounded (because $f\in C_b^2$), set $M:=\|f''\|_\infty$.
Then $|f''(x+\theta z)-f''(x)|\le 2M$, so
\[
\left|\frac12\int_{|z|>\delta}\!\!\bigl(f''(x+\theta z)-f''(x)\bigr)z^2\,
\frac{1}{\sqrt{2\pi h}}e^{-z^2/(2h)}\,dz\right|
\le
M\int_{|z|>\delta} z^2\,\frac{1}{\sqrt{2\pi h}}e^{-z^2/(2h)}\,dz.
\]
Write $z=\sqrt{h}\,w$ (so $dz=\sqrt{h}\,dw$). Then
\[
\int_{|z|>\delta} z^2\,\frac{1}{\sqrt{2\pi h}}e^{-z^2/(2h)}\,dz
=
h \int_{|w|>\delta/\sqrt{h}} w^2\,\frac{1}{\sqrt{2\pi}}e^{-w^2/2}\,dw
\le
h \int_{|w|>\delta/\sqrt{h}} w^2\,\phi(w)\,dw,
\]
where $\phi(w)=\frac{1}{\sqrt{2\pi}}e^{-w^2/2}$. As $h\downarrow0$, the threshold
$\delta/\sqrt{h}\to\infty$, and since $w^2\phi(w)$ is integrable over $\R$,
\[
\int_{|w|>\delta/\sqrt{h}} w^2\,\phi(w)\,dw \longrightarrow 0.
\]
Hence the contribution from $|z|>\delta$ is $o(h)$.

Combining (i) and (ii), we obtain $|R_h(x)|\le \frac{\varepsilon}{2}h + o(h)$, hence $R_h(x)=o(h)$.
Now divide \eqref{eq:main_term} by $h$ and let $h\downarrow0$:
\[
\lim_{h\downarrow0}\frac{T_h f(x)-f(x)}{h}
=
\frac12 f''(x).
\]
This is exactly $Gf(x)=\frac12 f''(x)$.
\end{proof}

\begin{remark}
    The fact that $G=\tfrac12\Delta$ reflects that Brownian motion spreads mass in space according to the heat equation.
\end{remark}

\paragraph{Why a generator?} Introducing the generator allows us to characterize Markov processes through local
properties, establish connections with partial differential equations, and formulate
martingale characterizations of the dynamics. The generator therefore serves as the
fundamental object that links probabilistic evolution, analytic structure, and
infinitesimal behavior.

\subsection{Martingale characterization}
To close this section, we want to describe how a Martingale arises naturally from a Brownian motion when we understand it as a Markov process. It captures the idea that, once
the infinitesimal drift prescribed by the generator is removed, the remaining evolution of
the process should exhibit no systematic trend. This property is expressed by identifying
certain functionals of the process that are martingales.

\paragraph{Why?} The Martingale characterization is fundamental because it describes the dynamics of the process
directly in terms of its sample paths, without reference to finite-time transition
probabilities. It forms the basis of the martingale problem approach, which is a central
tool for proving existence and uniqueness of Markov processes and for connecting
probabilistic dynamics with analytic structures such as partial differential equations.

\begin{proposition}[Proposition 3.2]
Let $f\in C_b^2(\R)$ and define
\[
M_t:=f(B_t)-f(B_0)-\int_0^t Gf(B_s)\,ds,
\qquad t\ge0.
\]
Then $(M_t)_{t\ge0}$ is a martingale with respect to $(\widehat{\mathcal F}_t)_{t\ge0}$.
\end{proposition}

\begin{proof}
Fix $0\le s<t$. We show that $\E[M_t\mid \widehat{\mathcal F}_s]=M_s$.

\medskip
\Step{1: conditional expectation of the increment $f(B_t)-f(B_s)$}
By the Markov property (time-homogeneous) and the definition of the semigroup,
\[
\E\!\left[f(B_t)\mid \widehat{\mathcal F}_s\right]
=
T_{t-s}f(B_s).
\]
Therefore,
\begin{equation}\label{eq:increment_f}
\E\!\left[f(B_t)-f(B_s)\mid \widehat{\mathcal F}_s\right]
=
T_{t-s}f(B_s)-f(B_s).
\end{equation}

\medskip
\Step{2: rewrite $T_{t-s}f(B_s)-f(B_s)$ as an integral of $T_u Gf(B_s)$}
For each fixed $x$, consider the function $u\mapsto T_u f(x)$.
Using the semigroup property $T_{u+h}f = T_u(T_h f)$, we can write the difference quotient:
\[
\frac{T_{u+h}f(x)-T_u f(x)}{h}
=
T_u\left(\frac{T_h f - f}{h}\right)(x).
\]
Letting $h\downarrow0$ and using that $f\in C_b^2(\R)$ so that
$\frac{T_h f - f}{h}\to Gf$ pointwise and is uniformly bounded for small $h$
(because $Gf=\tfrac12 f''$ is bounded), we may pass the limit through $T_u$ to obtain
\begin{equation}\label{eq:dTu}
\frac{d}{du}T_u f(x)=T_u Gf(x).
\end{equation}
Integrating \eqref{eq:dTu} from $u=0$ to $u=t-s$ gives
\begin{equation}\label{eq:duhamel}
T_{t-s}f(x)-f(x)=\int_0^{t-s} T_u Gf(x)\,du.
\end{equation}
Apply \eqref{eq:duhamel} with $x=B_s$:
\begin{equation}\label{eq:increment_as_integral}
T_{t-s}f(B_s)-f(B_s)=\int_0^{t-s} T_u Gf(B_s)\,du.
\end{equation}

\medskip
\Step{3: compute the conditional expectation of the drift term}
We claim that for each $r\ge s$,
\begin{equation}\label{eq:cond_drift}
\E\!\left[Gf(B_r)\mid \widehat{\mathcal F}_s\right]=T_{r-s}Gf(B_s).
\end{equation}
Indeed, this is the Markov property applied to the bounded measurable function $Gf$.

Now use Fubini's theorem (justified since $Gf$ is bounded and $t-s<\infty$) and \eqref{eq:cond_drift}:
\begin{align*}
\E\!\left[\int_s^t Gf(B_r)\,dr \,\Big|\, \widehat{\mathcal F}_s\right]
&=\int_s^t \E\!\left[Gf(B_r)\mid \widehat{\mathcal F}_s\right]\,dr\\
&=\int_s^t T_{r-s}Gf(B_s)\,dr
=\int_0^{t-s} T_u Gf(B_s)\,du.
\end{align*}
Comparing with \eqref{eq:increment_as_integral}, we obtain
\begin{equation}\label{eq:key_cancel}
\E\!\left[f(B_t)-f(B_s)\mid \widehat{\mathcal F}_s\right]
=
\E\!\left[\int_s^t Gf(B_r)\,dr \,\Big|\, \widehat{\mathcal F}_s\right].
\end{equation}

\medskip
\Step{4: conclude the martingale property}
Using \eqref{eq:key_cancel},
\begin{align*}
\E[M_t\mid \widehat{\mathcal F}_s]
&=\E\!\left[f(B_t)-f(B_0)-\int_0^t Gf(B_r)\,dr \,\Big|\, \widehat{\mathcal F}_s\right]\\
&=\E\!\left[f(B_s)-f(B_0)-\int_0^s Gf(B_r)\,dr \,\Big|\, \widehat{\mathcal F}_s\right]\\
&\quad+\E\!\left[f(B_t)-f(B_s)-\int_s^t Gf(B_r)\,dr \,\Big|\, \widehat{\mathcal F}_s\right]\\
&=M_s + 0.
\end{align*}
Hence $(M_t)_{t\ge0}$ is a martingale.
\end{proof}

\begin{remark}
Once the generator is known, the expression
\[
f(B_t)-f(B_0)-\int_0^t Gf(B_s)\,ds
\]
has zero conditional drift.
This is not accidental: it is the probabilistic counterpart of the fundamental theorem of calculus
applied along the random path.
The resulting martingale property is a powerful tool, as it allows us to connect analytic properties of $f$
to probabilistic properties of $B$.    
\end{remark}




\section{Strong Markov Property and its consequencies}

% --- Intuitive introduction to Section: Strong Markov Property, Reflection Principle,
%     and the Running Maximum ---

The results in this section deepen the Markov viewpoint developed earlier by showing that
Brownian motion can ``restart itself'' not only at deterministic times, but also at random
times determined by the path itself.
This strengthened form of memorylessness is the key that allows us to perform path
transformations, such as reflection at a hitting time, and to derive precise information
about extreme values of the process.

Conceptually, this is one of the most powerful features of Brownian motion:
local path symmetries, combined with memorylessness at stopping times, lead to exact global
distributional results such as the law of the running maximum.

Throughout this section, $(B_t)_{t\ge0}$ is a standard Brownian motion adapted to the
completed, right--continuous filtration $(\widehat{\mathcal F}_t)_{t\ge0}$.
We use the exponential martingales
\[
M_t(\theta):=\exp\!\left(i\theta B_t+\frac12\theta^2 t\right),
\qquad \theta\in\R,
\]
and the Optional Stopping Theorem (stated earlier).

\subsection{Strong Markov property}

The Markov property formalizes the idea that the future evolution of a stochastic process
depends on the present state but not on the past. In its basic form, this property is
formulated at fixed deterministic times: once the value of the process at time $t$ is
known, the conditional distribution of the future after time $t$ is independent of what
happened before $t$.

In many applications, however, the relevant time at which one wishes to ``restart'' the
process is not deterministic but random. Typical examples include the first time a
particle hits a boundary, the moment a process exceeds a given threshold, or the time at
which an observation is made based on the evolution of the process itself. Such random
times are described mathematically by stopping times.

The strong Markov property asserts that the Markovian behavior of the process persists
even when conditioning on such random times. Informally, it states that once a stopping
time $\tau$ occurs, the future evolution of the process after time $\tau$ behaves like a
fresh copy of the process started from the random state $X_\tau$, and is independent of
the past before $\tau$. In other words, the process ``forgets its history'' not only at
fixed times, but also at random times determined by its own evolution.

This property is essential for understanding the behavior of Markov processes under
random interventions and for analyzing events defined in terms of first hitting times,
exit times, or optimal stopping problems. From a conceptual point of view, the strong
Markov property reflects a deep self-similarity of Markov processes: whenever the process
is observed at a stopping time, the probabilistic laws governing its future evolution are
the same as those governing the original process, except for a shift in time and a change
in the initial condition.

The formal definition and results that follow make this intuition precise by combining
the Markov property with the notion of stopping times and conditional expectations. The
path--space framework developed earlier provides the natural setting in which this
property can be stated and proved rigorously.

\begin{theorem}[Strong Markov property]
Let $T$ be an a.s.\ finite stopping time. Define the shifted process
\[
B_t^{\,T}:=B_{T+t}-B_T,\qquad t\ge0.
\] Let $T$ be an a.s.\ finite stopping time. Then:
\begin{enumerate}
\item For each $t\ge0$, $B_t^{\,T}\sim\mathcal N(0,t)$.
\item For each $t\ge0$, $B_t^{\,T}$ is independent of $\widehat{\mathcal F}_T$.
\item The process $(B_t^{\,T})_{t\ge0}$ has independent, stationary increments and continuous paths.
Consequently, $(B_t^{\,T})_{t\ge0}$ is a standard Brownian motion, independent of $\widehat{\mathcal F}_T$.
\end{enumerate}
\end{theorem}

\begin{proof}
\Step{1: conditional characteristic function of a single increment}

Fix $t\ge0$ and $\theta\in\R$. For each $n\in\mathbb N$, set $T_n:=T\wedge n$ (bounded stopping time).
Using the martingale property of $(M_u(\theta))_{u\ge0}$ and the Optional Stopping Theorem at the bounded
stopping time $T_n$, we first note that $M_{u\wedge T_n}(\theta)$ is a martingale.
Hence for $t\ge0$,
\[
\E\!\left[M_{T_n+t}(\theta)\mid \widehat{\mathcal F}_{T_n}\right]=M_{T_n}(\theta)
\quad\text{a.s.}
\]
Write $M_{T_n+t}(\theta)=\exp(i\theta B_{T_n+t})\exp(\tfrac12\theta^2(T_n+t))$ and similarly for $M_{T_n}$.
Cancelling the common factor $\exp(i\theta B_{T_n})\exp(\tfrac12\theta^2 T_n)$ gives
\begin{equation}\label{eq:cond_char_Tn}
\E\!\left[\exp\!\bigl(i\theta(B_{T_n+t}-B_{T_n})\bigr)\,\Big|\,\widehat{\mathcal F}_{T_n}\right]
=\exp\!\left(-\frac12\theta^2 t\right).
\end{equation}
The right-hand side is deterministic and is the characteristic function of $\mathcal N(0,t)$.

\medskip
\Step{2: pass from $T_n$ to $T$}

Since $T_n\uparrow T$ a.s.\ and $B$ has continuous paths, we have
$B_{T_n}\to B_T$ and $B_{T_n+t}\to B_{T+t}$ a.s., hence
\[
\exp\!\bigl(i\theta(B_{T_n+t}-B_{T_n})\bigr)\to \exp\!\bigl(i\theta(B_{T+t}-B_T)\bigr)
\quad\text{a.s.}
\]
Moreover the integrand has modulus $1$, so dominated convergence applies.
Also $\widehat{\mathcal F}_{T_n}\uparrow \widehat{\mathcal F}_T$ (usual conditions).
Taking limits in \eqref{eq:cond_char_Tn} yields
\begin{equation}\label{eq:cond_char_T}
\E\!\left[\exp\!\bigl(i\theta(B_{T+t}-B_T)\bigr)\,\Big|\,\widehat{\mathcal F}_T\right]
=\exp\!\left(-\frac12\theta^2 t\right)\quad\text{a.s.}
\end{equation}

\medskip
\Step{3: conclude distribution and independence for one increment}

Equation \eqref{eq:cond_char_T} says: conditional on $\widehat{\mathcal F}_T$, the characteristic function of
$B_{T+t}-B_T$ is the same deterministic function $\exp(-\tfrac12\theta^2 t)$.
Therefore $B_{T+t}-B_T\sim \mathcal N(0,t)$.

Moreover, since the conditional characteristic function does not depend on $\widehat{\mathcal F}_T$,
we get independence: for any bounded $\widehat{\mathcal F}_T$--measurable random variable $H$,
\[
\E\!\left[H\,e^{i\theta(B_{T+t}-B_T)}\right]
=\E\!\left[H\,\E\!\left(e^{i\theta(B_{T+t}-B_T)}\mid \widehat{\mathcal F}_T\right)\right]
=\exp\!\left(-\frac12\theta^2 t\right)\E[H],
\]
which is exactly the factorization that characterizes independence.

\medskip
\Step{4: finite-dimensional distributions and independent increments}

Fix $0\le t_1<\dots<t_m$ and $\theta_1,\dots,\theta_m\in\R$.
Consider the random vector
\[
\bigl(B_{T+t_1}-B_T,\dots,B_{T+t_m}-B_T\bigr).
\]
It suffices to compute its conditional characteristic function given $\widehat{\mathcal F}_T$.
Using the tower property and repeatedly applying \eqref{eq:cond_char_T} to increments,
one obtains
\[
\E\!\left[\exp\!\left(i\sum_{k=1}^m \theta_k(B_{T+t_k}-B_T)\right)\Big|\widehat{\mathcal F}_T\right]
=
\exp\!\left(-\frac12\sum_{j=1}^m \Big(\sum_{k=j}^m\theta_k\Big)^2 (t_j-t_{j-1})\right),
\quad (t_0:=0),
\]
which is deterministic and coincides with the characteristic function of the corresponding Gaussian
vector for standard Brownian motion. In particular, this shows that $(B_t^{\,T})_{t\ge0}$
has stationary Gaussian increments, and by the same factorization argument as in Step 3,
these increments are independent of $\widehat{\mathcal F}_T$ and mutually independent.

\medskip
\Step{5: continuity}

For each $\omega$, $t\mapsto B_{T(\omega)+t}(\omega)-B_T(\omega)$ is continuous because $B(\omega)$ is continuous.
Thus $(B_t^{\,T})_{t\ge0}$ has continuous paths.

Combining Steps 4 and 5 proves that $(B_t^{\,T})_{t\ge0}$ is a standard Brownian motion and is independent of
$\widehat{\mathcal F}_T$.
\end{proof}

\medskip
\paragraph{Why exponential martingales?}
To prove the strong Markov property without stochastic calculus, the key tool is the family of
exponential martingales
\[
M_t(\theta)=\exp\!\left(i\theta B_t+\tfrac12\theta^2 t\right).
\]
These martingales encode the characteristic functions of Brownian increments.
Applying the Optional Stopping Theorem at bounded approximations of a stopping time allows us
to compute conditional characteristic functions of post--$T$ increments.
The crucial observation is that these conditional characteristic functions turn out to be
deterministic, which simultaneously identifies the Gaussian law of the increments and proves
their independence from the past.

\begin{remark}[From Markov families to the strong Markov property]
The notions of a Markov family, the (ordinary) Markov property, and the strong Markov
property all formalize, at different levels, the idea that a stochastic process can be
``restarted''. The distinction between them lies in the type of restart that is allowed.

A Markov family $\{\mathbb P_x\}_{x\in E}$ specifies the law of the process when it is
started from a given deterministic state at a fixed initial time. The (ordinary) Markov
property extends this idea by allowing the restart to occur at a fixed deterministic time,
but from a random state determined by the evolution of the process. Finally, the strong
Markov property asserts that the same restart principle remains valid when the restart
time itself is random and depends on the past of the process, provided it is a stopping
time.

\medskip
\noindent
These three levels of restartability can be summarized as follows:
\[
\begin{array}{c|c|c}
\text{Concept} & \text{Restart time} & \text{Restart state} \\ \hline
\text{Markov family} & \text{fixed (initial time)} & \text{deterministic} \\
\text{Markov property} & \text{fixed} & \text{random} \\
\text{Strong Markov property} & \text{random (stopping time)} & \text{random}
\end{array}
\]
\end{remark}


\subsection{Reflection principle}

One of the central questions when studying stochastic processes, and Brownian motion in
particular, is how the process behaves relative to barriers or thresholds. For instance,
one may ask for the probability that a Brownian path exceeds a given level before a fixed
time, or how the maximum of the process over a time interval is distributed. Such
questions are inherently path--dependent and cannot be answered solely by looking at the
distribution of the process at a single time.

The reflection principle provides a powerful and remarkably simple tool to address these
problems. Its key insight is that, due to the symmetry and independence of increments of
Brownian motion, certain path events can be transformed into simpler ones by reflecting
parts of the trajectory at an appropriate level. In particular, paths that cross a given
barrier can be put into one--to--one correspondence with paths that end at reflected
positions, without changing their probability.

Intuitively, the idea is the following. Consider a Brownian path that first hits a level
$a$ at some random time. From that hitting time onward, the future increments of the
process are independent of the past. By reflecting the path after this first hitting time
at the level $a$, one obtains a new path that has the same probability as the original one
but ends on the opposite side of the barrier. This symmetry allows probabilities involving
the running maximum of the process to be expressed in terms of probabilities involving
the final value of the process alone.

The reflection principle thus converts difficult questions about maxima and boundary
crossings into tractable distributional computations. It plays a fundamental role in the
derivation of hitting probabilities, distribution of the running maximum, and first--passage
times for Brownian motion. Conceptually, it relies on the same ideas that underlie the
strong Markov property: the ability to ``restart'' the process at a random hitting time
and to exploit the invariance of its probabilistic structure under time shifts and
reflections.

The formal statement and proof make this intuition precise by constructing an explicit
reflection mapping on path space and showing that it preserves the law of Brownian motion.

Fix $a\in\R$ and define the hitting time
\[
H_a:=\inf\{t>0:\,B_t=a\}.
\]
Define the reflected process $(\widetilde B_t)_{t\ge0}$ by
\[
\widetilde B_t:=
\begin{cases}
B_t, & t\le H_a,\\[1mm]
2a-B_t, & t>H_a.
\end{cases}
\]

\begin{theorem}[Reflection principle]
The process $(\widetilde B_t)_{t\ge0}$ is a standard Brownian motion.
\end{theorem}

\begin{proof}
We verify the defining properties via the strong Markov property at $H_a$.

\medskip
\Step{1: rewrite the post-hitting part using a fresh Brownian motion}

By the strong Markov property, the shifted process
\[
W_t:=B_{H_a+t}-B_{H_a},\qquad t\ge0,
\]
is a standard Brownian motion independent of $\widehat{\mathcal F}_{H_a}$.
On $\{t>H_a\}$ we have $B_t=B_{H_a}+W_{t-H_a}=a+W_{t-H_a}$, hence
\[
\widetilde B_t = 2a - B_t = 2a-(a+W_{t-H_a})=a-W_{t-H_a},\qquad t>H_a.
\]
Thus, pathwise,
\[
\widetilde B_t=
\begin{cases}
B_t, & t\le H_a,\\
a-W_{t-H_a}, & t>H_a.
\end{cases}
\]

\medskip
\Step{2: show the reflected increments after $H_a$ have Brownian law}

Since $W$ is Brownian, $-W$ is also Brownian (symmetry of the Gaussian law).
Therefore the process $t\mapsto a-W_t$ is a Brownian motion started at $a$.

Moreover, because $W$ is independent of $\widehat{\mathcal F}_{H_a}$, the post-$H_a$ evolution of
$\widetilde B$ is independent of the pre-$H_a$ history.

\medskip
\Step{3: finite-dimensional distribution}

Fix times $0\le t_1<\dots<t_m$. Condition on $H_a$ and distinguish whether each $t_k\le H_a$ or $t_k>H_a$.
On $\{H_a>t_m\}$ we have $\widetilde B_{t_k}=B_{t_k}$ for all $k$.
On $\{H_a\le t_1\}$ we have $\widetilde B_{t_k}=a-W_{t_k-H_a}$ and the vector has the same law as a Brownian
motion started at $a$ (by Step 2).
On mixed events (some times before, some after $H_a$), the joint law factorizes because the increments of $W$
are independent of $\widehat{\mathcal F}_{H_a}$, yielding exactly the Brownian finite-dimensional distributions.

\medskip
\Step{4: continuity}

$\widetilde B$ is continuous because $B$ is continuous and at $t=H_a$ we have $\widetilde B_{H_a}=B_{H_a}=a$,
so there is no jump at the reflection time.

Hence $(\widetilde B_t)_{t\ge0}$ is a standard Brownian motion.
\end{proof}

\subsection{Running maximum: distribution and a key joint identity}

In many applications, the behavior of a stochastic process is determined not only by its
value at a fixed time, but by the extreme values it attains over an interval. Examples
include the highest level reached by a fluctuating signal, the maximum displacement of a
particle up to a given time, or the maximal drawdown of a financial asset. Such quantities
depend on the entire past evolution of the process and are therefore intrinsically
path--dependent.

For a stochastic process $(B_t)_{t \ge 0}$, the running maximum
\[
S_t:=\sup_{0\le u\le t} B_u,\qquad t\ge0.
\]
summarizes this path--dependent information in a single random variable. Unlike $X_t$,
the running maximum is not Markovian by itself and cannot be described using only the
finite--dimensional distributions of the process at fixed times. Understanding its
distribution therefore requires tools that exploit the structure of the underlying
process, rather than just its one--time marginals.

In the case of Brownian motion, the running maximum is particularly tractable due to the
combination of symmetry, independent increments, and the strong Markov property. These
features make it possible to relate events involving $S_t$ to simpler events involving
the terminal value $B_t$ through the reflection principle. As a result, probabilities
concerning the maximum of Brownian motion can often be computed explicitly.

The next identity expresses a pathwise symmetry of Brownian motion with respect to the level
$a$ and will use it to obtain the distribution of the running maximum. On the event $\{S_t \ge a\}$, the process reaches the level $a$ at some (random) time
before $t$. By the strong Markov property, the future evolution after this first hitting
time is independent of the past and has the same law as a fresh Brownian motion started
from $a$. Reflecting the increments of the path after this hitting time across the level
$a$ produces a new path with the same probability, while mapping terminal values below $a$
to terminal values above $a$. As a consequence, events of the form $\{S_t \ge a,\; B_t \in
E\}$ with $E \subset (-\infty,a]$ -- Brownian paths that reach the
level $a$ at some time before $t$ but end below it -- are in one--to--one correspondence with events
$\{B_t \ge a,\; 2a-B_t \in E\}$ -- paths that end above $a$ whose reflected terminal
positions lie in the same region below $a$. The proposition therefore allows probabilities involving
the running maximum to be expressed in terms of the one--dimensional distribution of
$B_t$, and serves as the key tool for computing the law of $S_t$.

\begin{proposition}[Reflection identity]
Let $a\in\R$ and $t>0$. For any Borel set $E\subset(-\infty,a]$,
\[
\Pp\bigl(S_t\ge a,\; B_t\in E\bigr)=\Pp\bigl(2a-B_t\in E,\; B_t\ge a\bigr).
\]
In particular, for $y\ge0$,
\[
\Pp\bigl(S_t\ge a,\; B_t\le a-y\bigr)=\Pp(B_t\ge a+y).
\]
\end{proposition}

\begin{proof}
Observe that \[
\{S_t\ge a\}=\{H_a\le t\}.
\]Consider the path transformation $\omega\mapsto \widetilde\omega$ defined by reflection at $H_a(\omega)$,
i.e.\ $\widetilde B$ as above. By the reflection principle theorem, $\widetilde B$ is itself a Brownian motion,
so it has the same law as $B$.

Now note the event mapping:
on $\{H_a\le t\}$ (equivalently $\{S_t\ge a\}$), we have $\widetilde B_t=2a-B_t$;
on $\{H_a>t\}$, we have $\widetilde B_t=B_t$ and also $S_t<a$.
Therefore, for any Borel $E\subset(-\infty,a]$,
\begin{align*}
\Pp\bigl(S_t\ge a,\; B_t\in E\bigr)
&=\Pp\bigl(H_a\le t,\; B_t\in E\bigr)\\
&=\Pp\bigl(H_a\le t,\; 2a-\widetilde B_t\in E\bigr)
\qquad (\text{since }\widetilde B_t=2a-B_t \text{ on }\{H_a\le t\})\\
&=\Pp\bigl(H_a\le t,\; \widetilde B_t\in 2a-E\bigr),
\end{align*}
where $2a-E:=\{2a-x:\,x\in E\}\subset[a,\infty)$.
But on $\{\widetilde B_t\in 2a-E\}\subset\{\widetilde B_t\ge a\}$, we necessarily have $H_a\le t$,
so the event $\{H_a\le t\}$ may be dropped:
\[
\Pp\bigl(H_a\le t,\; \widetilde B_t\in 2a-E\bigr)=\Pp\bigl(\widetilde B_t\in 2a-E\bigr).
\]
Finally, since $\widetilde B$ is Brownian, $\widetilde B_t\overset{d}=B_t$, hence
\[
\Pp\bigl(\widetilde B_t\in 2a-E\bigr)=\Pp\bigl(B_t\in 2a-E\bigr)
=\Pp\bigl(2a-B_t\in E\bigr).
\]
Restricting to $E\subset(-\infty,a]$ forces $2a-B_t\le a$, i.e.\ $B_t\ge a$, so
\[
\Pp\bigl(2a-B_t\in E\bigr)=\Pp\bigl(2a-B_t\in E,\; B_t\ge a\bigr),
\]
which proves the first identity. The second follows by choosing $E=(-\infty,a-y]$.
\end{proof}

We can now compute the running maximum from the distribution of $B_t$.

\begin{corollary}[Distribution of the maximum]
For $a\ge0$ and $t>0$,
\[
\Pp(S_t\ge a)=2\,\Pp(B_t\ge a)
=2\int_{a}^{\infty}\frac{1}{\sqrt{2\pi t}}e^{-x^2/(2t)}\,dx
=2\int_{a/\sqrt t}^{\infty}\frac{1}{\sqrt{2\pi}}e^{-u^2/2}\,du.
\]
\end{corollary}

\begin{proof}
Take $y=0$ in the particular case of the reflection identity:
\[
\Pp(S_t\ge a,\; B_t\le a)=\Pp(B_t\ge a).
\]
Also, since $\{S_t\ge a\}$ implies $B_t\in(-\infty,a]\cup(a,\infty)$ and the boundary has probability $0$,
we have the disjoint decomposition
\[
\Pp(S_t\ge a)=\Pp(S_t\ge a,\; B_t\le a)+\Pp(S_t\ge a,\; B_t>a).
\]
But $\{B_t>a\}\subset\{S_t\ge a\}$, so the second term is $\Pp(B_t>a)=\Pp(B_t\ge a)$.
Thus $\Pp(S_t\ge a)=2\,\Pp(B_t\ge a)$.
Finally, $B_t\sim \mathcal N(0,t)$ gives the Gaussian integral formulas, with the change of variables
$x=\sqrt t\,u$ in the last equality.
\end{proof}


The study of the running maximum thus serves two complementary purposes. On the one hand,
it provides insight into the extremal behavior of Brownian paths and leads to concrete
distributional results. On the other hand, it illustrates how path--dependent quantities
can be analyzed by combining stopping times, the strong Markov property, and symmetry
arguments. The results of this section play a central role in the analysis of hitting
times, barrier crossing problems, and related questions in stochastic process theory.